\section{Related Work}
\label{sec:related-work}

\textbf{Float-to-string algorithms.}
Burger and Dybvig~\cite{burger1996printing} established the
shortest-representation framework using multiprecision arithmetic.
Loitsch~\cite{loitsch2010grisu} introduced Grisu, the first
integer-only approach, handling $\sim$99.5\% of inputs without fallback.
Adams~\cite{adams2018ryu} proposed Ry\={u}, achieving full coverage with
fixed-width 128-bit arithmetic and lookup tables.
Jeon~\cite{jeon2021dragonbox} developed Dragonbox with improved table
compression.  Giulietti~\cite{giulietti2022schubfach} presented
\schubfach, which our implementation follows.  Champagne Gareau and
Lemire~\cite{champagnegareau2026} provide a comprehensive survey of
these algorithms.

All prior work focuses on the algorithms in isolation.  Our contribution
is the first comparison \emph{inside a specification-constrained pipeline}
where the formatting stage and surrounding canonicalization logic are
non-trivial.

\textbf{JSON canonicalization.}
\rfcjcs~\cite{rfc8785} is the primary standard.  Prior to RFC~8785,
ad-hoc canonicalization schemes existed in various ecosystems (e.g.,
JSON-LD canonicalization, OLPC canonical JSON).  None specify the
floating-point formatting requirement with the precision of
\ecma~\cite{ecma262}.

\textbf{Benchmark methodology.}
Our conformance-before-performance approach draws on the principle that
correctness is a prerequisite for meaningful performance comparison, as
advocated in systems benchmarking methodology literature.  The use of
permutation tests and bootstrap confidence intervals follows modern
statistical practice for performance evaluation.
